% !TEX root = ../thesis.tex

\chapter{Introduction} \label{chp:intro}

People love quantum computing I'm focused on an aspect of it.

Background is stuff that I don't want to explain to people in the later chapters. Superconducting basics. Dressed basis two qubit gates, stabilizers for error correcting codes, basic circuit quantization using the devoret recipe.

Explaining dispersive readout, voltage control on transmon, flux tuning, etc. The basics. Transmission lines.

Discuss protection vs. error correction.

Noise estimates.

\section{Quantum Computation}


\section{Circuit Quantum Electrodynamics (Circuit QED)}

\subsection{Superconductivity}

\subsection{LC Resonator}

\subsection{Transmission Lines}

\subsection{Basic Circuit Quantization}

\section{Cooper Pair Box \& Transmon}

\subsection{Tunable Transmon}

\subsection{Single Qubit Gates}

\subsection{Two Qubit Gates}

\subsection{Dispersive Readout}

\section{Fluxonium}

\subsection{Single Qubit Gates}

\subsection{Two Qubit Gates}

\subsection{Dispersive Readout}

\section{Protected Qubits}

\subsection{Estimating Lifetimes}

\subsection{Cos $2\varphi$ Qubit}

\subsection{$0-\pi$ Qubit}

\section{Quantum Error Correction (QEC)}

\subsection{Stabilizer Formalism}





% Here is an example of using the macros
% \verb2\singlespacing2 and \verb2\doublespacing2:

% \singlespacing	% <------------------------------

% This paragraph was preceded by the
% command \verb2\singlespacing2.
% See the Specifications of the Grad School for instructions
% about when single spacing is appropriate in a thesis.

% \doublespacing	% <------------------------------

% And now, here is an example of using the macros
% \verb2\begin{singlespace}2 and \verb2\end{singlespace}2;
% another way to get single-spacing:

% \begin{singlespace}	% <-----------------------------
% Two cases are studied in the present work which differ only in the
% boundary conditions.  Each different boundary condition model a
% different source of instability.  The boundary of the first case
% consists of a steady, axisymmetric sidewall radial velocity boundary
% and a time-dependent, non-axisymmetric endwall axial velocity
% boundary.  The second case is studied with a fixed impermeable axial
% velocity along the endwall and a combination axisymmetric steady and
% non-axisymmetric unsteady radial velocity along the sidewall.
% \end{singlespace}	% <-----------------------------
